\section{Ejecución del bloque III: P09--P12 y correcciones asignadas}
\label{sec:bloque3}
\textbf{Fecha: 18 de octubre de 2026.} Se implementó el control local de fallos,
evidencia y presupuestos. P10 permanece abierto para calibración humana y P12
para medir utilidad musical real. No se afirma superioridad de relevancia.

\subsection{Reubicación y corrección de E14--E16}
E14 pertenece al bloque II/P07: la identidad del plan debe distinguir cambios
efectivos de una permutación de listas. La huella ahora canoniza descriptores y
registros de restricciones sin borrar procedencia. La caché de recuperación usa
la misma independencia de orden. Una prueba del grafo rechaza candidatos 1 y 2
una sola vez, conserva la versión 1 y termina por ventana agotada; antes repetía
el candidato 1 tres veces y creaba versiones 2 y 3. Un descriptor efectivamente
distinto sigue cambiando la identidad. Esto no demuestra equivalencia semántica
entre textos distintos: E13 continúa pendiente.

E15 pertenece al bloque II/P05/P08: la salida de aclaración ahora verifica JSON
finito igual que un plan listo. Números NaN/infinito y objetos opacos se rechazan;
el orquestador conserva su fallback controlado. Se probaron ambas rutas y la
serialización del estado final recuperado. Se conserva el rechazo explícito del
dato inválido, sin sustituirlo por una confianza inventada.

E16 pertenece al bloque I/P01/P02: los mensajes de router, workers y juez usan
una salida que escapa símbolos incompatibles con la codificación de consola.
Un diagnóstico deja de interrumpir una inicialización correcta. UTF-8 sigue
siendo una opción operativa; en CP1252 algunos símbolos se ven escapados.
Ambas codificaciones completaron las mismas 96 pruebas.

\subsection{P09: recuperación parcial y continuidad}
\textbf{Antes.} Una excepción de una rama abortaba la recuperación completa;
un refinador inválido detenía el flujo. \textbf{Ahora.} Cada rama informa éxito,
vacío, ausencia de solicitud o fallo. Los resultados útiles de la otra rama se
conservan; sus pesos efectivos se renormalizan y se exponen junto a los pesos
solicitados y la política. Las restricciones originales no se eliminan.
Si falla el refinador, se conserva el plan válido y se deshabilitan nuevos
intentos del refinador para esa solicitud. Si falla toda una nueva recuperación,
se restaura el plan, versión y ranking obtenidos juntos en la última búsqueda
válida, aplicando las exclusiones de esa versión.

\textbf{Evidencia y contrapartida.} Pruebas del grafo y HTTP cubren fallos
parciales, fallos de ambas ramas y restauración. Una petición mixta sigue
informando sonido sin evaluar aunque la letra se apruebe. Esta política mejora
continuidad, pero puede reducir cobertura y cambiar scores; no garantiza la
misma calidad que disponer de ambas ramas. P13 debe explicar contribuciones.

\subsection{P10: evidencia lírica, abstención y reintento acotado}
\textbf{Antes.} El juez veía los primeros 3000 caracteres y emitía texto o una
herramienta de rechazo. \textbf{Ahora.} Se solicita JSON con estado, motivo y
citas literales. Las citas se contrastan con fragmentos realmente suministrados.
Las letras de hasta 12000 caracteres se incluyen completas; las más largas
aportan comienzo, centro y final dentro de ese presupuesto. La cobertura parcial
impide aprobar el conjunto. La abstención no se transforma en rechazo semántico
ni activa el refinador; puede continuar la revisión de otros candidatos.

Un formato o evidencia inválidos permiten un solo reintento correctivo;
una excepción del proveedor no se reintenta por esta capa. El cliente del juez
desactiva sus reintentos internos y registra intentos de modelo. Se conserva
compatibilidad con veredictos antiguos, etiquetados como \texttt{legacy}, sin
atribuirles citas verificadas. Esta fue la decisión inicial del bloque;
la sección~\ref{sec:correcciones_pendientes} documenta su retirada posterior.
Los demás candidatos siguen sin evaluar.

\textbf{Evidencia y contrapartida.} Una prueba encuentra y cita evidencia
después del carácter 4000; otras cubren citas inventadas, cobertura incompleta,
abstención, errores y el máximo de dos invocaciones. El contexto puede ser mayor
que antes y el máximo de salida pasó de 256 a 1024 tokens; ello puede aumentar
coste y latencia. Una cita exacta no prueba que el juicio semántico sea correcto.
\texttt{evaluation/judge\_metrics.py} calcula falsos aprobados/rechazados,
abstención y error solo con etiquetas humanas. Sin pares anotados devuelve
métricas nulas: la calibración y E13 siguen pendientes, no se inventaron etiquetas.

\subsection{P11: presupuestos independientes}
\textbf{Antes.} Tres visitas a candidatos podían agotarse sin llamar al juez
si faltaban metadatos; se buscaba una ventana fija. \textbf{Ahora.} Por solicitud
hay límites separados: 20 visitas, 3 evaluaciones del juez (hasta 6 intentos de
modelo por el reintento), 2 refinamientos y 1 ampliación de recuperación de
top-10 a top-20 por rama. Las omisiones de metadatos no consumen evaluaciones
del juez. Los contadores y el motivo de parada se exponen en la API.

La ampliación solo se intenta al agotar una ventana saturada; se filtran las
exclusiones previas. Agotar esa ventana no demuestra ausencia de canciones
válidas en todo el catálogo; la API no declara agotamiento global.
\textbf{Pruebas.} Cuatro metadatos ausentes antes de un candidato válido,
presupuesto total de visitas, ampliación que encuentra un candidato nuevo y
agotamiento de la ampliación. Revisar más candidatos puede consumir más trabajo;
los límites son configurables al construir el grafo y no una estimación monetaria
ni un límite global de todas las solicitudes HTTP del proveedor.

\subsection{P12: comparación pareada y política de refinamiento}
\textbf{Antes.} El refinador se invocaba tras rechazos sin evidencia de mejora
musical. \textbf{Ahora.} Está desactivado por defecto en la aplicación; la opción
\texttt{SPELLFIND\_ENABLE\_REFINEMENT=true} permite uso experimental. Solo se
considera tras rechazo lírico, con contrato válido y presupuesto disponible.
No se activa por abstención, metadatos ausentes ni recuperación acústica sola.

Se capturaron 36 ejecuciones sin refinamiento y 36 con él, pareadas por consulta
y repetición, con hashes iguales de fuentes, tests, datos y locks. Ambas series
simuladas tuvieron 12 aprobaciones líricas, 6 validaciones parciales,
12 recuperaciones acústicas y 6 aclaraciones, sin errores. El tratamiento añadió
18 invocaciones al refinador y 18 búsquedas líricas; no cambió el número de
evaluaciones del juez. Los veredictos y rankings son sintéticos: no demuestran
mejor o peor relevancia. Los tiempos simulados y su dispersión se conservan en
el informe, sin convertirlos en latencia de modelos o precio del proveedor.

\texttt{evaluation/compare.py} rechaza snapshots distintos y muestra diferencias
pareadas. La evaluación real escribió estado bloqueado por falta de
\texttt{GROQ\_API\_KEY}; las etiquetas humanas siguen vacías. Mantener la opción
desactivada limita trabajo adicional no justificado, pero renuncia a posibles
beneficios todavía no medidos. Para decidir su activación rutinaria se requieren
comparaciones reales de relevancia, violaciones negativas, abstención y coste,
con anotaciones y condiciones comparables. P12 no queda cerrado por la simulación.

\subsection{Verificación, compatibilidad y pendientes}
\textbf{Archivos.} E14/E15 cambian \texttt{intent.py} y la clave de caché en
\texttt{app.py}; E16 añade \texttt{console.py} y lo integra en router, workers
y juez. P09/P11 usan \texttt{workflow.py}, \texttt{retrieval.py} y \texttt{app.py}.
P10 cambia \texttt{judge.py}; P12 actualiza aplicación, ejemplo de configuración
y \texttt{evaluation/run.py}, y añade comparación y métricas de calibración.
La interfaz reconoce abstención. \texttt{tests/test\_block3.py} añade regresiones;
se actualizan expectativas anteriores para la continuidad y presupuestos nuevos.

\texttt{scripts/verify\_block3.py} ejecutó 96 pruebas con salida CP1252 y 96 con
UTF-8, ambas series correctas. La evidencia está en \texttt{validation/block3/}.
No se instalaron dependencias, descargaron pesos ni hicieron llamadas al modelo.
Los informes de bloques anteriores y el diagnóstico original conservan su
carácter histórico. El documento se actualizó en fuente; no se certifica en esta
entrega una compilación PDF.

Los estados y campos añadidos requieren que clientes reconozcan abstenciones,
degradaciones y presupuestos. El nombre histórico \texttt{judge\_iterations}
sigue contando evaluaciones; los intentos de modelo están en contadores separados.
P13--P16 son el siguiente bloque de implementación. P10/P12 y la fidelidad
semántica E13 mantienen sus verificaciones abiertas; no se renumeran ni se
consideran resueltas al avanzar.
